\documentclass[11pt,a4paper]{article}
\usepackage[margin=25mm]{geometry}
\usepackage{newtxtext}
\usepackage{newtxmath}
\usepackage{microtype}
\usepackage{bm}
\usepackage{mathtools}
\usepackage[numbers,sort&compress]{natbib}
\usepackage[colorlinks,allcolors=blue]{hyperref}

\newcommand{\Tr}{\operatorname{Tr}}
\newcommand{\diag}{\operatorname{diag}}
\newcommand{\ee}{\mathrm{e}}
\newcommand{\ii}{\mathrm{i}}
\newcommand{\trans}{\mathsf{T}}
\newcommand{\mc}[1]{\mathcal{#1}}
\newcommand{\bb}[1]{\bm{#1}}
\allowdisplaybreaks

\title{Split-orthogonal positivity in bosonic AFQMC\\
    \large Physical annihilation space and Nambu reduction}
\author{}
\date{\today}

\begin{document}
\maketitle

The split-orthogonal determinant theorem gives a sufficient condition for
sign-free number-conserving fermionic auxiliary-field quantum Monte Carlo
(AFQMC)~\cite{wang2015prl}. Its bosonic application depends on the operator
space on which the matrix acts. This note derives the bosonic determinant
sign, explains the grand-canonical obstruction when the group acts on
physical annihilation operators, and shows why the same group works in the
bosonic $\{+,+,+\}$ time-reversal class. A stable two-mode thermal state makes
the distinction explicit. We then compare the number-conserving
time-reversal symmetry (TRS) and Gaussian reflection-positivity (RP) conditions.

\section{The determinant theorem and its bosonic counterpart}

For a fixed auxiliary-field configuration, write a number-conserving Gaussian
product as
\begin{equation}
    \hat U=\prod_{\ell=1}^{L}\ee^{\hat h_\ell},
    \qquad
    \hat h_\ell=\hat{\bb a}^{+}X_\ell\hat{\bb a},
    \qquad
    \mc B=\prod_{\ell=1}^{L}\ee^{X_\ell}.
    \label{eq:number_product}
\end{equation}
Here $\hat{\bb a}$ contains all physical annihilation operators and
$\hat{\bb a}^{+}$ is the corresponding creation row. Real scalar terms
in the slices supply a positive overall factor and are omitted for now.
The fermionic trace and the bosonic determinant representative are
\begin{equation}
    W_{\rm F}=\det(I+\mc B),
    \qquad
    W_{\rm B}^{\rm det}=\det(I-\mc B)^{-1}.
    \label{eq:weights}
\end{equation}
The second expression equals the ordinary bosonic Fock trace in its
convergent domain. Outside that domain it is a meromorphic continuation.

Let $\eta=\diag(I_n,-I_n)$. A real matrix belongs to $O(n,n)$ when
$M^\trans\eta M=\eta$. Its four connected components are labeled by the
signs of the determinants of the two diagonal $n\times n$ blocks.
The identity component is $O^{++}(n,n)$. Taking determinants gives
$(\det M)^2=1$, and continuity from the identity therefore fixes
\begin{equation}
    \det M=1
    \qquad\text{for every }M\in O^{++}(n,n).
    \label{eq:unit_det}
\end{equation}
This is a property of the group, independent of particle statistics.
Real generators satisfying
$X_\ell^\trans\eta+\eta X_\ell=0$ produce an ordered product in this component.

The split-orthogonal theorem states that $\det(I+M)$ is nonnegative on
$O^{++}(n,n)$ and nonpositive on $O^{--}(n,n)$~\cite{wang2015prl}.
For $M\in O^{++}(n,n)$, the matrix $-M$ lies in $O^{++}$ when $n$ is even
and in $O^{--}$ when $n$ is odd. Applying the theorem to $-M$ yields
\begin{equation}
    \boxed{(-1)^n\det(I-M)\geq0,\qquad M\in O^{++}(n,n).}
    \label{eq:minus_sign}
\end{equation}
Thus a direct substitution $\mc B=M$ in Eq.~\eqref{eq:weights} gives a
positive bosonic determinant representative for $n$ even and a negative one
for $n$ odd, away from poles. For example, the elementary boost
\begin{equation}
    R(x)=
    \begin{pmatrix}\cosh x&\sinh x\\ \sinh x&\cosh x\end{pmatrix}
    \in O^{++}(1,1),
    \qquad
    \det[I-R(x)]=-4\sinh^2(x/2)
    \label{eq:boost}
\end{equation}
already exhibits the odd-$n$ sign.

\section{The obstruction in physical annihilation space}
\label{sec:annihilation_obstruction}

The convergence of a bosonic occupation sum requires every physical
annihilation mode to decay. In particular, an ordinary trace-class Gaussian
product requires $\rho(\mc B)<1$, where $\rho$ is the spectral radius.
For a general nonnormal matrix, this necessary spectral condition alone
does not ensure trace-classness. Equation~\eqref{eq:unit_det} implies
$\rho(\mc B)\geq1$ whenever $\mc B\in O^{++}(n,n)$, so this group cannot
describe a convergent Gaussian trace on the full physical annihilation
space. The positive sign in Eq.~\eqref{eq:minus_sign} for even $n$ refers
to the determinant continuation.

When every slice preserves the same $\eta$, the obstruction also has a
many-body interpretation. Bosonic commutation relations give
\begin{equation}
    \hat Q^{+}=\frac12\hat{\bb a}^{+}
                       \eta(\hat{\bb a}^{+})^\trans,
    \qquad
    [\hat h_\ell,\hat Q^{+}]
    =\frac12\hat{\bb a}^{+}
      (X_\ell\eta+\eta X_\ell^\trans)
      (\hat{\bb a}^{+})^\trans=0.
    \label{eq:pair_tower}
\end{equation}
The final equality follows from the Lie-algebra condition and $\eta^2=I$.
Each number-conserving slice annihilates the vacuum, so every product in
Eq.~\eqref{eq:number_product} acts with the same eigenvalue on
\begin{equation}
    |0\rangle,\quad
    \hat Q^{+}|0\rangle,\quad
    (\hat Q^{+})^2|0\rangle,\quad\ldots .
    \label{eq:tower_states}
\end{equation}
After normalization these are infinitely many orthogonal states in
different particle-number sectors. Restoring scalar slice terms multiplies
all of them by the same factor. If an exact HS representation imposes
this common group condition on every configuration, the full thermal
operator inherits the same tower as its vacuum. A physical Hermitian
Hamiltonian with a finite vacuum energy then has a divergent grand-canonical
partition function. The symmetric pair in Eq.~\eqref{eq:pair_tower} is
nonzero for bosons; the analogous symmetric creation bilinear vanishes
for fermions.

A chemical potential changes the number-conserving propagation to
$z\mc B$, with fugacity $z=\ee^{\beta\mu}$. It also changes the group relation:
\begin{equation}
    (z\mc B)^\trans\eta(z\mc B)=z^2\eta.
    \label{eq:fugacity}
\end{equation}
Consequently, the theorem at $z=1$ does not determine the sign of
$\det(I-z\mc B)$. Projection onto a fixed total particle number removes
the infinite tower from the trace, but its weight is the coefficient
$[z^{N_{\rm p}}]\det(I-z\mc B)^{-1}$, whose sign requires a separate argument.
For instance, $-I_4\in O^{++}(2,2)$ has a positive inverse determinant
at $z=1$, while its one-particle trace is $-4$.

\section{Why Nambu reduction admits a stable split-orthogonal propagation}

Nambu space contains both annihilation and creation operators. Even a
number-conserving problem has the doubled propagation
\begin{equation}
    \hat\chi=(\hat{\bb a},(\hat{\bb a}^{+})^\trans)^\trans,
    \qquad
    \mc M=
    \begin{pmatrix}\mc B&0\\0&\mc B^{-\trans}\end{pmatrix},
    \qquad
    \det\mc M=1.
    \label{eq:nambu_product}
\end{equation}
A contracting annihilation eigenvalue is accompanied by an expanding
creation eigenvalue. This reciprocal structure is compatible with a
finite thermal trace. It records the canonical algebra of the doubled
operator space.

In the bosonic $\{+,+,+\}$ class, the time reversals are physical,
pairwise anticommuting antiunitary maps that each square to $+I$.
Their product in pairs selects a half-dimensional Nambu sector.
The reduced propagation, denoted $M_{\rm red}$, preserves a real symmetric
form with signature $(n,n)$ and lies in $O^{++}(n,n)$ after a basis change.
The selected operators have a commutator matrix with $n$ positive and
$n$ negative directions. A negative direction corresponds to selecting
a creation operator in place of its annihilation partner.

This choice changes the trace factor. At a stable diagonal reference point,
an annihilation direction with propagation $\ee^s$ contributes
$\ee^{s/2}/(1-\ee^s)$ to the centered Gaussian trace. Selecting the creation
direction instead gives
\begin{equation}
    \frac{\ee^{-s/2}}{1-\ee^{-s}}
    =-\frac{\ee^{s/2}}{1-\ee^s}.
    \label{eq:reversed_direction}
\end{equation}
There is one minus sign per negative commutator direction. The Gaussian
trace formula carries this factor to mixed directions and ordered
products.\footnote{The sector construction and the general reduced
Gaussian trace are derived in the companion note
\emph{Time-Reversal Symmetries for Sign-Problem-Free Quantum Monte Carlo}.}
Since the reduced generators are traceless, the result in this class is
\begin{equation}
    \boxed{
    W_{\rm B}=\ee^\Xi
       \frac{(-1)^n}{\det(I-M_{\rm red})}>0
    }
    \label{eq:reduced_weight}
\end{equation}
where $\Xi$ is the sum of the real scalar terms separated from the
centered quadratic generators. Positivity follows from
Eq.~\eqref{eq:minus_sign}, away from poles, with the branch fixed by
the operator trace. Both $\det\mc M=1$ and $\det M_{\rm red}=1$ hold.
The physical content of the reduced operators supplies the sign factor.

\subsection*{A two-mode thermal example}

Take two bosonic modes with
\begin{equation}
    H=\varepsilon(n_b+n_c),\qquad
    x=\beta\varepsilon>0,\qquad q=\ee^{-x}<1.
    \label{eq:thermal_model}
\end{equation}
The physical annihilation propagation is $\mc B=qI_2$, and
$Z=(1-q)^{-2}$ is finite. The antiunitary maps
$K$, $\ii\sigma_xK$, and $\ii\sigma_zK$ on $(b,c)^\trans$
preserve $H$, anticommute pairwise, and square to $+I$.
For the $+\ii$ sector of the product of the first two maps, choose
\begin{equation}
    d_1=\frac{b-c}{\sqrt2},\qquad
    d_2=\frac{b^{+}+c^{+}}{\sqrt2},
    \qquad
    [d_\mu,d_\nu^{+}]=
       \begin{pmatrix}1&0\\0&-1\end{pmatrix}_{\mu\nu}.
    \label{eq:thermal_sector}
\end{equation}
With the propagation convention
$\ee^{-\hat h}\bb d\,\ee^{\hat h}=M_{\rm red}\bb d$ and
$\hat h=-\beta H$, one obtains
\begin{equation}
    M_{\rm red}=\diag(q,q^{-1}),\qquad
    M_{\rm red}^\trans
       \begin{pmatrix}0&1\\1&0\end{pmatrix}M_{\rm red}
       =\begin{pmatrix}0&1\\1&0\end{pmatrix}.
    \label{eq:thermal_metric}
\end{equation}
An orthogonal change of basis diagonalizes this metric to
$\diag(1,-1)$ and transforms $M_{\rm red}$ into $R(-x)$ of
Eq.~\eqref{eq:boost}. Its determinant is one and its eigenvalue $q^{-1}$
is larger than one, precisely because $d_2$ is a creation direction.

Centering the two quadratic modes gives
$\hat h=-x(n_b+n_c+1)+x$, hence $\Xi=x$. Direct evaluation of
Eq.~\eqref{eq:reduced_weight} now gives
\begin{equation}
    \det(I-M_{\rm red})=-\frac{(1-q)^2}{q},
    \qquad
    Z=\ee^x\frac{-1}{\det(I-M_{\rm red})}
      =\frac1{(1-q)^2}.
    \label{eq:thermal_weight}
\end{equation}
The same boost that causes an occupation-growth problem when assigned
to two physical annihilation modes describes a stable thermal state
when assigned to one annihilation and one creation direction.

\section{Relation to number-conserving TRS and Gaussian RP}

Kramers TRS on physical annihilation space pairs complex eigenvalues by
conjugation and forces real eigenvalues to have even multiplicity.
It therefore protects $\det(I-\mc B)>0$ away from poles, and real
chemical-potential shifts preserve this condition. Split orthogonality
instead enforces reciprocal spectral pairing. The two matrix conditions
are not nested. Two unequal boost blocks $R(x)\oplus R(y)$, with
$x,y>0$ and $x\ne y$, have four distinct real eigenvalues and no Kramers
symmetry. Two identical blocks give an algebraic overlap with
conjugate-copy TRS and RP. Conversely, the stable thermal propagation
$qI_2$ has Kramers TRS and RP but has determinant $q^2<1$.
The overlap involving boosts is a statement about matrix conditions;
the grand-canonical obstruction of Sec.~\ref{sec:annihilation_obstruction} still applies when those
boosts act on physical annihilation modes.

The number-conserving restriction simplifies the current Gaussian RP
condition. Consider a fixed division into independent modes $b$ and $c$,
with an antilinear multiplicative reflection $\theta(b)=c$ and
$\theta(b^{+})=c^{+}$. Site phases can be absorbed into the $c$ modes.
Write its Gaussian generator condition as
\begin{equation}
    \hat h=\hat q+\theta(\hat q)
       +\sum_{\mu\nu}C_{\mu\nu}X_\mu\theta(X_\nu)+\xi,
    \qquad
    C=C^\dagger\succeq0,\quad \xi\in\mathbb R,
    \quad
    \bb X=(b_1,\ldots,b_m,b_1^{+},\ldots,b_m^{+})^\trans .
    \label{eq:rp_condition}
\end{equation}
Conservation of $N_b+N_c$ sets the two diagonal blocks of $C$ to zero,
since they multiply $bc$ and $b^{+}c^{+}$. Positive semidefiniteness
then forces its off-diagonal blocks to vanish as well, by
$|C_{\mu\nu}|^2\le C_{\mu\mu}C_{\nu\nu}$. Thus $C=0$, and
\begin{equation}
    \hat h_\ell=
       \hat{\bb b}^{+}D_\ell\hat{\bb b}
       +\hat{\bb c}^{+}\overline{D_\ell}\hat{\bb c}+\xi_\ell .
    \label{eq:rp_number}
\end{equation}
This is the conjugate-copy limit of the Kramers block form
\begin{equation}
    X_\ell=
    \begin{pmatrix}D_\ell&F_\ell\\
    -\overline{F_\ell}&\overline{D_\ell}\end{pmatrix}.
    \label{eq:kramers_block}
\end{equation}
For these number-preserving reflections and Gaussian slice conditions,
RP is therefore contained in Kramers TRS.

The inclusion can be strict for a fixed HS representation. Define
$S_\alpha=(b^{+},c^{+})\sigma_\alpha(b,c)^\trans$ and take the stable
number-conserving model
\begin{equation}
    H=\epsilon(N_b+N_c)+g_xS_x^2+g_zS_z^2,
    \qquad \epsilon,g_x,g_z>0.
    \label{eq:trs_example}
\end{equation}
Gaussian HS decoupling generates imaginary potentials
$\ii\phi_x\sigma_x$ and $\ii\phi_z\sigma_z$. Both preserve
$\ii\sigma_yK$. Their nonzero commutator prevents simultaneous
diagonalization into the two one-mode blocks of Eq.~\eqref{eq:rp_number}.
This HS representation satisfies Kramers TRS beyond the stated Gaussian
RP condition. A different decomposition of the same Hamiltonian would
require its own comparison.

Finally, if all three $\{+,+,+\}$ maps preserve the physical annihilation
space, their product is itself an antiunitary symmetry with square $-I$.
They then imply the number-conserving Kramers criterion. The additional
role of this class appears when pairing requires the full Nambu
description. There, Eq.~\eqref{eq:reduced_weight} combines the
split-orthogonal determinant sign with the signs of the selected
commutator directions. This representation-dependent mechanism is
compatible with a finite bosonic thermal trace.

\bibliographystyle{unsrt}
\bibliography{references}
\end{document}
